\par\Needspace{6\baselineskip}\medskip
\hypertarget{native-127-article094}{}\pdfbookmark[2]{94.}{article-094}
\begin{center}94.\end{center}
Tractatio aequationum 12, $12^{\ast}$ in artt. 91, 92, 93 explicata innixa est suppositioni, angulum $g$ non esse nimis magnum, certe infra limitem $66^0\,25'$, ultra quem tabulam III non extendimus. Quoties haec suppositio locum non habet, aequationes illae tantis artificiis non indigent: poterunt enim \textit{forma non mutata} tutissime semper ac commodissime tentando solvi. \textit{Tuto} scilicet, quoniam valor expressionis $\dfrac{2g-\sin 2g}{\sin g^3}\mathpunct{\raisebox{.45ex}{,}}$ in qua $2g$ in partibus radii exprimendum esse sponte patet, pro valoribus maioribus ipsius $g$ \textit{omni praecisione} computari potest per tabulas trigonometricas, quod utique fieri nequit, quamdiu $g$ est angulus parvus: \textit{commode}, quoniam loci heliocentrici tanto intervallo ab invicem distantes vix unquam ad determinationem orbitae penitus adhuc incognitae adhibebuntur, ex orbitae autem cognitione qualicunque valor approximatus ipsius $g$ nullo propemodum negotio per aequationem 1 vel 3 art. 88 demanat: denique e valore approximato ipsius $g$ valor correctus, aequationi 12 vel $12^{\ast}$ omni quae desideratur praecisione satisfaciens, semper paucis tentaminibus eruetur. Ceterum quoties duo loci heliocentrici propositi plus una revolutione integra complectuntur, memorem esse oportet, quod ab anomalia excentrica totidem revolutiones completae absolutae erunt, ita ut anguli $E'-E$, $v'-v$ vel ambo \SourceLeaf{128}iaceant inter 0 et $360^0$, vel ambo inter multipla similia totius peripheriae, adeoque $f$ et $g$ vel simul inter 0 et $180^0$, vel inter multipla similia semiperipheriae. Quodsi tandem orbita omnino incognita esset, neque adeo constaret, utrum corpus coeleste, transeundo a radio vectore primo ad secundum, descripserit partem tantum revolutionis, an insuper revolutionem integram unam seu plures, problema nostrum nonnunquam plures solutiones diversas admitteret: attamen huic casui in praxi vix unquam occursuro hic non immoramur.
